\subsection{Relación de beneficio / costo}

\subsubsection{Identificación de la Pérdida Económica}

La Tabla~\ref{tab:perdidas_consolidadas} sintetiza el impacto financiero originado por las causas raíz evaluadas en el levantamiento inicial, 
considerando el tiempo dedicado por el personal contable a corregir inconsistencias y los costos por minuto asociados a dichas tareas.

\par
\noindent
\begin{minipage}{\linewidth}
    \tituloTablaAPA{tab:perdidas_consolidadas}{Pérdida económica mensual y anual estimada por causa raíz identificada}

    \begingroup
    \renewcommand{\baselinestretch}{1}
    \footnotesize
    \setlength{\tabcolsep}{4pt}
    \renewcommand{\arraystretch}{1.25}

    \begin{tabularx}{\linewidth}{@{}p{1.4cm}X p{3.4cm}r r@{}}
        \toprule
        \textbf{Código} & \textbf{Descripción de la Causa Raíz} & \textbf{Tipo de Pérdida} & \textbf{Mensual (S/)} & \textbf{Anual (S/)} \\ \midrule
        CA1 & Información laboral almacenada en archivos separados & Conciliación manual de legajos & S/ 186.60 & S/ 2,239.20 \\ \addlinespace
        CA2 & Circulación de versiones de planilla sin control & Revisión de copias desactualizadas & S/ 75.20 & S/ 902.40 \\ \addlinespace
        CA3 & Transcripción manual de asistencias e incidencias & Corrección de omisiones y tareo & S/ 111.75 & S/ 1,341.00 \\ \addlinespace
        CA4 & Modificación manual de fórmulas y parámetros & Revisión de diferencias de cálculo & S/ 110.70 & S/ 1,328.40 \\ \addlinespace
        CA7 & Preparación separada de archivos de pago y PLAME & Ajustes en transferencias y SUNAT & S/ 74.85 & S/ 898.20 \\ \addlinespace
        CA8 & Generación y distribución individual de boletas & Tiempo de envío manual y reenvíos & S/ 62.48 & S/ 749.76 \\ \midrule
        \textbf{Total} & & & \textbf{S/ 621.58} & \textbf{S/ 7,458.96} \\ \bottomrule
    \end{tabularx}
    \endgroup

    \notaTabla{Pérdida mensual calculada sobre los minutos de corrección y el costo horario del personal. Pérdida anual computada a 12 meses. Elaboración propia.}
\end{minipage}
\finBloqueAPA

La pérdida anual de \textbf{S/ 7,458.96} representa el costo de ineficiencia operativa medido estrictamente en horas de 
retrabajo del personal contable durante el levantamiento del proceso. No incluye el impacto 
económico potencial de sanciones o multas inspectivas que el sistema previene de manera sistemática.

\clearpage
\subsubsection{Identificación de los Beneficios del Sistema}

Al estandarizar los flujos de trabajo e incorporar reglas de validación automáticas en el software, se establece una tasa conservadora 
del \textbf{85.0\,\% de eficiencia de recuperación} sobre las pérdidas operativas detectadas (Tabla~\ref{tab:eficiencia_recuperacion}).

\par
\noindent
\begin{minipage}{\linewidth}
    \tituloTablaAPA{tab:eficiencia_recuperacion}{Evaluación de la eficiencia de recuperación de pérdidas mediante el sistema web}

    \begingroup
    \renewcommand{\baselinestretch}{1}
    \footnotesize
    \setlength{\tabcolsep}{4pt}
    \renewcommand{\arraystretch}{1.25}

    \begin{tabularx}{\linewidth}{@{}r c r c@{}}
        \toprule
        \textbf{Pérdidas Anuales (S/)} & \textbf{Representan \%} & \textbf{Costo de Inversión (S/)} & \textbf{\% Eficiencia de Recuperación} \\ \midrule
        S/ 7,458.96 & 100.0\,\% & S/ 60,906.79 & 85.0\,\% \\ \bottomrule
    \end{tabularx}
    \endgroup

    \notaTabla{Porcentaje de mitigación operativa esperado tras la automatización de los procesos. Elaboración propia.}
\end{minipage}
\finBloqueAPA

Aplicando dicho factor del 85.0\,\%, se determina el ahorro económico anual directo generado por el aplicativo (Tabla~\ref{tab:total_ahorro_beneficio}):

\begin{equation}
    \text{Ahorro Anual (Beneficio Operativo)} = \text{Pérdidas Anuales} \times \text{Eficiencia de Recuperación}
\end{equation}

\begin{equation*}
    \text{Ahorro Anual (Beneficio Operativo)} = \text{S/ } 7,458.96 \times 0.85 = \text{\textbf{S/ 6,340.12}}
\end{equation*}

\par
\noindent
\begin{minipage}{\linewidth}
    \tituloTablaAPA{tab:total_ahorro_beneficio}{Cálculo del beneficio económico anual neto por ahorro operativo}

    \begingroup
    \renewcommand{\baselinestretch}{1}
    \footnotesize
    \setlength{\tabcolsep}{4pt}
    \renewcommand{\arraystretch}{1.25}

    \begin{tabularx}{\linewidth}{@{}X r@{}}
        \toprule
        \textbf{Concepto} & \textbf{Monto (S/)} \\ \midrule
        Pérdidas anuales en el proceso manual previo & S/ 7,458.96 \\ 
        Sobrecosto residual proyectado tras la mejora (15\,\%) & S/ 1,118.84 \\ \midrule
        \textbf{TOTAL AHORRO ANUAL (BENEFICIO OPERATIVO)} & \textbf{S/ 6,340.12} \\ \bottomrule
    \end{tabularx}
    \endgroup

    \notaTabla{Beneficio económico derivado de la reducción de tiempos muertos y reprocesos contables. Elaboración propia.}
\end{minipage}
\finBloqueAPA

\subsubsection{Relación Beneficio / Costo (B/C)}

Dado que el desarrollo fue concebido para enriquecer el portafolio comercial de J\&P Periféricos S.A.C., 
el análisis financiero evalúa el beneficio global del proyecto: integra el ahorro operativo generado en las operaciones de prueba (S/ 6,340.12) 
y los ingresos anuales proyectados por la suscripción al servicio SaaS de 20 clientes empresariales ($20 \times \text{S/ } 280.00 \times 12 = \text{S/ } 67,200.00$), 
alcanzando un beneficio anual consolidado de \textbf{S/ 73,540.12}.

Contrastando este beneficio frente al costo total de la inversión calculado en el Capítulo~5 (S/ 60,906.79):

\begin{equation}
    \text{Relación } B/C = \frac{\text{Beneficio Anual Total}}{\text{Costo Total de Inversión}}
\end{equation}

\begin{equation*}
    \text{Relación } B/C = \frac{\text{S/ } 73,540.12}{\text{S/ } 60,906.79} = \mathbf{1.21}
\end{equation*}

\noindent
\textbf{Interpretación:} Al ser la relación Beneficio/Costo mayor a la unidad ($B/C = 1.21 > 1.00$), 
se demuestra que el proyecto es \textbf{financieramente viable y rentable}. Por cada sol (S/ 1.00) 
invertido en el desarrollo e implantación del sistema, se recupera el capital inicial y se genera un retorno neto adicional de S/ 0.21.

\subsubsection{Periodo de Recuperación de la Inversión}

El tiempo necesario para amortizar la totalidad de los recursos invertidos se calcula mediante la siguiente relación:

\begin{equation}
    \text{Periodo de Recuperación} = \frac{\text{Costo Total de Inversión}}{\text{Beneficio Anual Total}}
\end{equation}

\begin{equation*}
    \text{Periodo de Recuperación} = \frac{\text{S/ } 60,906.79}{\text{S/ } 73,540.12} = 0.8282\text{ años}
\end{equation*}

Transformando el resultado a meses y días:

\begin{equation*}
    0.8282\text{ años} \times 12\text{ meses/año} = 9.94\text{ meses} \longrightarrow \mathbf{9\text{ meses}}
\end{equation*}

\begin{equation*}
    0.94\text{ meses} \times 30\text{ días/mes} = 28.2\text{ días} \longrightarrow \mathbf{28\text{ días}}
\end{equation*}

\noindent
\textbf{Interpretación:} La inversión realizada se recupera íntegramente en un periodo de \textbf{9 meses y 28 días}. 
Este plazo confirma que el proyecto se amortiza en menos de un año fiscal, 
respaldando la solidez de la propuesta tecnológica y garantizando su sostenibilidad en el tiempo.